% Generated draft: requires mathematical and source review.
\documentclass[11pt]{article}
\usepackage[margin=25mm]{geometry}
\usepackage{fontspec}
\setmainfont{Latin Modern Roman}
\usepackage{amsmath,amssymb,mathtools,unicode-math}
\setmathfont{Latin Modern Math}
\usepackage{enumitem,needspace,xurl,hyperref}
\setlength{\parindent}{0pt}
\setlength{\parskip}{5pt}
\setlength{\emergencystretch}{3em}
\newcommand{\N}{\mathbb{N}}
\newcommand{\Z}{\mathbb{Z}}
\newcommand{\Q}{\mathbb{Q}}
\newcommand{\R}{\mathbb{R}}
\newcommand{\C}{\mathbb{C}}
\newcommand{\F}{\mathbb{F}}
\newcommand{\A}{\mathbb{A}}
\newcommand{\PP}{\mathbb{P}}
\newcommand{\E}{\mathbb{E}}
\newcommand{\eps}{\varepsilon}
\newcommand{\dd}{\,\mathrm{d}}
\begin{document}
\section*{UnsolvedMath ID: 700955}
\textbf{Source problem code:} catalogue entry
\subsection{Definitions and mathematical statement}
Let \(p\) be a prime. A pro-\(p\) group is a topological group isomorphic to an inverse limit of finite \(p\)-groups. Such a group \(G\) is torsion-free if \(g^n=1\), for \(g\in G\) and an integer \(n\geq 1\), implies \(g=1\). An open subgroup \(U\leq G\) carries the subspace topology; its index in \(G\) is a finite power of \(p\).

A pro-\(p\) presentation with one defining relation has the form
\[
H\cong F/\overline{\langle\!\langle r\rangle\!\rangle}_F,
\]
where \(F\) is a free pro-\(p\) group, \(r\in F\), the denominator is the smallest closed normal subgroup of \(F\) containing \(r\), and the isomorphism is topological. Write \(\operatorname{OneRel}_p(H)\) for the source's property that \(H\) is a pro-\(p\) group with a single defining relation.

The question attributed to O. V. Mel'nikov is whether, for every prime \(p\), every torsion-free pro-\(p\) group \(G\), and every open subgroup \(U\leq G\),
\[
\operatorname{OneRel}_p(U)\quad\Longrightarrow\quad
\operatorname{OneRel}_p(G).
\]
The supplied record does not specify finite generation, the rank of the free pro-\(p\) group in a presentation, or whether the presentation must be minimal. It also does not explicitly state conventions concerning trivial or redundant relators. These qualifications remain unresolved: the displayed quotient describes a one-relation presentation, but is not intended to settle the exact convention for \(\operatorname{OneRel}_p\). No finite-generation hypothesis has been added.
\subsection{Short English statement}
If an open subgroup of a torsion-free pro-\(p\) group has a single defining relation in the pro-\(p\) sense, must the whole group also have a single defining relation? The source leaves presentation conventions and finite-generation assumptions unspecified.
\subsection{Sources}
\begin{itemize}
\item Local UnsolvedMath registry: permanent ID \(700955\), exact source problem code \texttt{catalogue entry}, title \emph{Kourovka Notebook Problem 13.37}; supplied upstream revision \texttt{d371ee917ac9647fbb06cd6fc529deec6f3dda82}.
\item Supplied attribution: O. V. Mel'nikov, proposed in 1995; \emph{The Kourovka Notebook (Unsolved Problems in Group Theory)}, 13th issue (1995), Problem 13.37. The record also identifies the 21st issue edition (2026), \texttt{arXiv:1401.0300v48}. Its collection description names E. I. Khukhro and V. D. Mazurov as editors.
\item The supplied registry labels the problem \emph{open} and records an update date of 2026-10-16. This is supplied registry metadata, not an independently verified assessment of the literature. No external source was visited for this draft.
\end{itemize}
\end{document}
